% Proposed addition, 2026-08-26 meeting item: "case study, Budget Guard vs agent
% self-reporting, for a Table 1 task" (about half a page). NOT yet in the paper.
% Every number below is produced by scripts/make_case_study.py and stored in
% tables/case_study_sokoban.csv, tables/case_study_sokoban_tally.csv and
% tables/case_study_sokoban_gate.csv (Claude Opus 4.7 Sokoban run, seed-0 test split). Selection rule, fixed in the script before
% any rollout was inspected: the first rollout in the seed-0 test stream (rollout 7).
% The 0.51 / 0.92 figures are the verified Figure 2 values.
%
% Placement suggestion: end of Section 5.1, before "5.2 Robustness under deployment
% shift". Requires \usepackage{booktabs} and \usepackage{amssymb} (for \checkmark).
% \ref{tab:token_abort} must be replaced by the source's actual label for Table 3
% (token-abort economics); check with grep before pasting.

\paragraph{A worked example.}
Table~\ref{tab:case_study} follows the first held-out Sokoban rollout in the seed-0
test stream (Claude Opus 4.7, 2,500-token budget) turn by turn; the rollout is fixed
by that rule, not chosen by hand. At each prefix the agent reports its own
remaining-token interval and the guard reports a split-CQR interval computed from
the same prefix; the true remaining spend is known only in hindsight. The agent's
intervals are 500 to 800 tokens wide and cover only at the first prefix: from turn 2
they sit entirely above the truth, and the final self-report reserves 400 to 900
tokens for a run that needed 267. The guard's intervals are wider but contain the
truth at every prefix, and they support the decision a deployer actually has to make:
spend so far plus the guard's upper bound stays under the budget at every turn
(2,139, 2,178 and 2,282 tokens against 2,500), so the run could be certified to fit
from its first turn. The feasibility gate, calibrated at $\delta{=}0.10$ on the same
split, scores the run at 0.99, 0.95 and 0.99 against an abort threshold of 0.02, so
it says continue at every turn; the run finished inside budget and succeeded.
Because this is the first rollout of the stream, the online windowed-CQR intervals
coincide with the offline ones here. The example is representative rather than selected: across all 30
test rollouts of this split, split-CQR covers a larger fraction of prefixes than the
self-report on 22, ties on 6 and trails on 2 (windowed-CQR: 24, 5 and 1), consistent
with the Sokoban column of Figure~\ref{fig:coverage} (0.51 for the self-report
against 0.92 for the guard). On the same split the gate aborts none of the 16
feasible test runs and recovers 4.9\% of doomed-run spend; Table~\ref{tab:token_abort}
reports the five-seed means.

\begin{table}[t]
\centering
\caption{The first held-out Sokoban rollout in the seed-0 test stream, turn by turn
(Claude Opus 4.7, 2,500-token budget). \checkmark marks an interval that contains
the true remaining spend. The last column is spend so far plus the guard's upper
bound, which a deployer compares against the budget; the gate column is the
feasibility score and the decision at $\delta{=}0.10$ (abort below 0.02).}
\label{tab:case_study}
\footnotesize
\setlength{\tabcolsep}{4pt}
\begin{tabular}{rrrlllr}
\toprule
Turn & Spent & True rem. & Agent self-report & Guard interval (split-CQR) & Spent + upper & Gate \\
\midrule
1 & 557 & 1,090 & [1,000, 1,800] \checkmark & [460, 1,582] \checkmark & 2,139 $\leq$ 2,500 & 0.99 continue \\
2 & 1,076 & 571 & [600, 1,200] $\times$ & [210, 1,102] \checkmark & 2,178 $\leq$ 2,500 & 0.95 continue \\
3 & 1,380 & 267 & [400, 900] $\times$ & [162, 902] \checkmark & 2,282 $\leq$ 2,500 & 0.99 continue \\
\bottomrule
\end{tabular}
\end{table}

% Optional figure (meeting: "if space allows, a figure at the beginning to show how
% the framework processed it"). Same rollout as the table; figures/case_study_sokoban.png.
% If used, the table above can drop to the appendix, or vice versa; both is redundant.
\begin{figure}[t]
\centering
\includegraphics[width=0.7\linewidth]{figures/case_study_sokoban.png}
\caption{The same Sokoban rollout as Table~\ref{tab:case_study}. Grey bars are the
agent's self-reported remaining-token interval at each turn, green bars the guard's
split-CQR interval, black dots the true remaining spend, and the dashed red line the
budget still available. The self-report drifts above the truth after turn 1; the
guard's interval tracks it and stays under the remaining budget throughout, and the
gate (scores under each turn) says continue at every step.}
\label{fig:case_study}
\end{figure}

% Two sentences for Section 6 (Discussion and limitations), if the example goes in:
% "The worked example in Section 5.1 also shows what the guard does not do: its
% intervals are roughly twice as wide as the agent's, and the price of a calibrated
% answer is a less precise one. Where a deployer only needs a yes-or-no on whether
% the run still fits, that width costs nothing; where a tight point estimate
% matters, the interval is a bound, not a forecast."
